%! Author = Profesor
%! Date = 06/06/2024


\section{Ejercicios y Ejercicios EVAU}

\subsection{Ejercicios}


\begin{exercise}
    En medio ácido, el permanganato de potasio oxida el ion cloruro a cloro reduciéndose a manganeso (III). Escribe y ajusta la reacción usando el método del ion-electrón.
    \begin{solu}

        \textbf{Semirreacción de reducción} (\ch{Mn^{+7} -> Mn^{+3}}, gana 4 electrones):

        \begin{equation*}
            \ch{MnO4- + 8 H+ + 4 e- -> Mn^{3+} + 4 H2O}
        \end{equation*}

        \textbf{Semirreacción de oxidación} (\ch{Cl^{-1} -> Cl^{0}}, pierde 1 electrón por átomo):

        \begin{equation*}
            \ch{2 Cl- -> Cl2 + 2 e-}
        \end{equation*}

        Para igualar los electrones intercambiados (mínimo común múltiplo de 4 y 2 es 4), multiplicamos la semirreacción de oxidación por 2:

        \begin{equation*}
            \ch{4 Cl- -> 2 Cl2 + 4 e-}
        \end{equation*}

        Sumando ambas semirreacciones:

        \begin{equation*}
            \ch{MnO4- + 8 H+ + 4 Cl- -> Mn^{3+} + 4 H2O + 2 Cl2}
        \end{equation*}

        Pasando a la forma molecular con \ch{KMnO4} y \ch{HCl} (que aporta tanto los \ch{H+} como los \ch{Cl-} necesarios, incluyendo los 3 \ch{Cl-} que acompañan al \ch{Mn^{3+}} en el \ch{MnCl3} y 1 \ch{Cl-} para el \ch{KCl}):

        \begin{equation*}
            \ch{KMnO4 + 8 HCl -> MnCl3 + KCl + 2 Cl2 + 4 H2O}
        \end{equation*}

        Comprobación: K (1=1), Mn (1=1), O (4=4), H (8=8), Cl (\(8 = 3+1+4\)).

    \end{solu}
\end{exercise}

\begin{exercise}
    Las siguientes reacciones se dan en medio ácido:

    \begin{enumerate}
        \item \ch{MnO2\sld{} + Cl- \aq{} -> Mn\pch[2]\aq{} + Cl2\gas{}}
        \item \ch{I- \aq{} + Br2\lqd{} -> IO3- \aq{} + Br- \aq{}}
    \end{enumerate}
    \begin{solu}

        \textbf{1) \ch{MnO2 + Cl- -> Mn^2+ + Cl2}}

        Reducción (\ch{Mn^{+4} -> Mn^{+2}}, gana 2 e\textsuperscript{-}):

        \begin{equation*}
            \ch{MnO2 + 4 H+ + 2 e- -> Mn^{2+} + 2 H2O}
        \end{equation*}

        Oxidación (\ch{Cl^{-1} -> Cl^{0}}, pierde 2 e\textsuperscript{-} por cada \ch{Cl2}):

        \begin{equation*}
            \ch{2 Cl- -> Cl2 + 2 e-}
        \end{equation*}

        Como ambas ceden/ganan 2 electrones, se suman directamente:

        \begin{equation*}
            \ch{MnO2 + 4 H+ + 2 Cl- -> Mn^{2+} + 2 H2O + Cl2}
        \end{equation*}

        En forma molecular (con \ch{HCl}):

        \begin{equation*}
            \ch{MnO2 + 4 HCl -> MnCl2 + Cl2 + 2 H2O}
        \end{equation*}

        \textbf{2) \ch{I- + Br2 -> IO3- + Br-}}

        Oxidación (\ch{I^{-1} -> I^{+5}}, pierde 6 e\textsuperscript{-}):

        \begin{equation*}
            \ch{I- + 3 H2O -> IO3- + 6 H+ + 6 e-}
        \end{equation*}

        Reducción (\ch{Br^{0} -> Br^{-1}}, gana 2 e\textsuperscript{-} por cada \ch{Br2}):

        \begin{equation*}
            \ch{Br2 + 2 e- -> 2 Br-}
        \end{equation*}

        Igualando electrones (m.c.m. de 6 y 2 es 6), multiplicamos la reducción por 3:

        \begin{equation*}
            \ch{3 Br2 + 6 e- -> 6 Br-}
        \end{equation*}

        Sumando:

        \begin{equation*}
            \ch{I- + 3 H2O + 3 Br2 -> IO3- + 6 H+ + 6 Br-}
        \end{equation*}

        En forma molecular (con \ch{HI} y \ch{HBr}):

        \begin{equation*}
            \ch{HI + 3 H2O + 3 Br2 -> HIO3 + 6 HBr}
        \end{equation*}

    \end{solu}
\end{exercise}

\begin{exercise}
    Ajusta las siguientes reacciones por el método del ion-electrón:

    \begin{enumerate}[label=\letraminuscula]
        \item \ch{Cl2 + NaOH -> NaCl + NaClO + H2O}
        \item \ch{HNO3 + Cu -> Cu(NO3)2 + NO2 + H2O}
        \item \ch{HCl + K2Cr2O7 -> KCl + CrCl3 + Cl2 + H2O}
        \item \ch{KNO3 + Al + KOH + H2O -> NH3 + KAlO2}
        \item \ch{Na2SO4 + C -> Na2S + CO2}
        \item \ch{MnO4- + SO2 -> Mn\pch[2] + HSO4-}
        \item \ch{MnO2- + I- -> MnO4- + IO3-}
        \item \ch{I2 + HNO3 -> HIO3 + NO2 + H2O}
        \item \ch{MnO4- + SO3\mch[2] -> Mn\pch[2] + SO4\mch[2]} (medio ácido)
        \item \ch{TiO2 + Al -> Al2O3 + Ti}
        \item \ch{HNO3\aq{} + C\sld{} -> CO2\gas{} + NO2\gas{} + H2O\lqd{}}
        \item \ch{Bi(OH)3\aq{} + Na2SnO2\aq{} -> Bi\sld{} + Na2SnO3\aq{} + H2O\lqd{}}
    \end{enumerate}
    \begin{solu}

        \textbf{a) \ch{Cl2 + NaOH -> NaCl + NaClO + H2O}} (dismutación del cloro en medio básico)

        Reducción: \(\ch{Cl2 + 2 e- -> 2 Cl-}\). Oxidación: \(\ch{Cl2 + 4 OH- -> 2 ClO- + 2 H2O + 2 e-}\).

        Sumando (mismos electrones): \(\ch{2 Cl2 + 4 OH- -> 2 Cl- + 2 ClO- + 2 H2O}\), que simplificado da:

        \begin{equation*}
            \ch{Cl2 + 2 NaOH -> NaCl + NaClO + H2O}
        \end{equation*}

        \textbf{b) \ch{HNO3 + Cu -> Cu(NO3)2 + NO2 + H2O}} (ácido nítrico concentrado)

        Reducción: \(\ch{NO3- + 2 H+ + e- -> NO2 + H2O}\) (\(\times 2\)). Oxidación: \(\ch{Cu -> Cu^{2+} + 2 e-}\).

        \begin{equation*}
            \ch{Cu + 4 HNO3 -> Cu(NO3)2 + 2 NO2 + 2 H2O}
        \end{equation*}

        \textbf{c) \ch{HCl + K2Cr2O7 -> KCl + CrCl3 + Cl2 + H2O}}

        Reducción: \(\ch{Cr2O7^2- + 14 H+ + 6 e- -> 2 Cr^{3+} + 7 H2O}\). Oxidación: \(\ch{2 Cl- -> Cl2 + 2 e-}\) (\(\times 3\)).

        \begin{equation*}
            \ch{K2Cr2O7 + 14 HCl -> 2 KCl + 2 CrCl3 + 3 Cl2 + 7 H2O}
        \end{equation*}

        \textbf{d) \ch{KNO3 + Al + KOH + H2O -> NH3 + KAlO2}} (medio básico)

        Reducción: \(\ch{NO3- + 6 H2O + 8 e- -> NH3 + 9 OH-}\) (\(\times 3\)). Oxidación: \(\ch{Al + 4 OH- -> AlO2- + 2 H2O + 3 e-}\) (\(\times 8\)).

        Sumando y simplificando \ch{H2O} y \ch{OH-} sobrantes:

        \begin{equation*}
            \ch{3 KNO3 + 8 Al + 5 KOH + 2 H2O -> 3 NH3 + 8 KAlO2}
        \end{equation*}

        \textbf{e) \ch{Na2SO4 + C -> Na2S + CO2}}

        \ch{S}: \(+6 \to -2\) (gana 8 e\textsuperscript{-}); \ch{C}: \(0 \to +4\) (pierde 4 e\textsuperscript{-}). Se necesitan 2 mol de C por cada mol de S:

        \begin{equation*}
            \ch{Na2SO4 + 2 C -> Na2S + 2 CO2}
        \end{equation*}

        \textbf{f) \ch{MnO4- + SO2 -> Mn^2+ + HSO4-}} (medio ácido)

        Reducción: \(\ch{MnO4- + 8 H+ + 5 e- -> Mn^{2+} + 4 H2O}\) (\(\times 2\)). Oxidación: \(\ch{SO2 + 2 H2O -> HSO4- + 3 H+ + 2 e-}\) (\(\times 5\)).

        \begin{equation*}
            \ch{2 MnO4- + 5 SO2 + 2 H2O + H+ -> 2 Mn^{2+} + 5 HSO4-}
        \end{equation*}

        \textbf{g) \ch{MnO2- + I- -> MnO4- + IO3-}}

        \textit{Nota:} tal como está escrita, esta ecuación no se puede ajustar como proceso redox, ya que tanto el manganeso (que pasaría de un estado de oxidación \(+3\) en \ch{MnO2-} a \(+7\) en \ch{MnO4-}) como el yodo (que pasaría de \(-1\) en \ch{I-} a \(+5\) en \ch{IO3-}) se \textbf{oxidarían}, sin que exista ninguna especie que se reduzca. Es muy probable que se trate de una errata en el enunciado (por ejemplo, el reactivo de manganeso debería ser una especie donde el Mn se reduzca, o falte algún otro reactivo oxidante). Por ello, esta reacción no puede ajustarse tal como está planteada.

        \textbf{h) \ch{I2 + HNO3 -> HIO3 + NO2 + H2O}}

        Oxidación: \(\ch{I2 + 6 H2O -> 2 IO3- + 12 H+ + 10 e-}\). Reducción: \(\ch{NO3- + 2 H+ + e- -> NO2 + H2O}\) (\(\times 10\)).

        \begin{equation*}
            \ch{I2 + 10 HNO3 -> 2 HIO3 + 10 NO2 + 4 H2O}
        \end{equation*}

        \textbf{i) \ch{MnO4- + SO3^2- -> Mn^2+ + SO4^2-}} (medio ácido)

        Reducción: \(\ch{MnO4- + 8 H+ + 5 e- -> Mn^{2+} + 4 H2O}\) (\(\times 2\)). Oxidación: \(\ch{SO3^2- + H2O -> SO4^2- + 2 H+ + 2 e-}\) (\(\times 5\)).

        \begin{equation*}
            \ch{2 MnO4- + 5 SO3^2- + 6 H+ -> 2 Mn^{2+} + 5 SO4^2- + 3 H2O}
        \end{equation*}

        \textbf{j) \ch{TiO2 + Al -> Al2O3 + Ti}}

        \ch{Ti}: \(+4 \to 0\) (gana 4 e\textsuperscript{-}); \ch{Al}: \(0 \to +3\) (pierde 3 e\textsuperscript{-}). m.c.m.(4,3) = 12:

        \begin{equation*}
            \ch{3 TiO2 + 4 Al -> 2 Al2O3 + 3 Ti}
        \end{equation*}

        \textbf{k) \ch{HNO3 + C -> CO2 + NO2 + H2O}}

        Oxidación: \(\ch{C + 2 H2O -> CO2 + 4 H+ + 4 e-}\). Reducción: \(\ch{NO3- + 2 H+ + e- -> NO2 + H2O}\) (\(\times 4\)).

        \begin{equation*}
            \ch{C + 4 HNO3 -> CO2 + 4 NO2 + 2 H2O}
        \end{equation*}

        \textbf{l) \ch{Bi(OH)3 + Na2SnO2 -> Bi + Na2SnO3 + H2O}} (medio básico)

        Reducción: \(\ch{Bi(OH)3 + 3 e- -> Bi + 3 OH-}\) (\(\times 2\)). Oxidación: \(\ch{SnO2^2- + 2 OH- -> SnO3^2- + H2O + 2 e-}\) (\(\times 3\)).

        Sumando y cancelando los \ch{OH-} que se repiten en ambos lados:

        \begin{equation*}
            \ch{2 Bi(OH)3 + 3 Na2SnO2 -> 2 Bi + 3 Na2SnO3 + 3 H2O}
        \end{equation*}

    \end{solu}
\end{exercise}

%* Libro Santillana página 300
\begin{exercise}
    Se hace pasar una corriente eléctrica de \SI{5}{\ampere} durante \num{2.5} horas a través de una celda electrolítica que contiene una disolución de cloruro de cobre (II). Calcule:
    \begin{enumerate}[label=\letraminuscula]
        \item La masa de cobre metálico depositado en el cátodo.
        \item El volumen de cloro molecular medido en condiciones normales que se genera en el ánodo.
    \end{enumerate}
    \datos{} \(F = \num{96500}\ \dfrac{\text{C}}{\text{mol}}\); masa atómica Cu = \num{63.5}
    \begin{solu}

        La carga total que atraviesa la celda es:

        \begin{equation*}
            Q = I \cdot t = 5 \ \si{\ampere} \times (2,5 \times 3600) \ \si{\s} = \SI{45000}{\coulomb}
        \end{equation*}

        Los moles de electrones intercambiados son:

        \begin{equation*}
            n(e^-) = \dfrac{Q}{F} = \dfrac{45000}{96500} = 0,4663 \ \text{mol}
        \end{equation*}

        \textbf{a) Masa de cobre depositada en el cátodo}

        En el cátodo tiene lugar la reducción \(\ch{Cu^{2+} + 2 e- -> Cu}\), por lo que:

        \begin{equation*}
            n(\ch{Cu}) = \dfrac{n(e^-)}{2} = \dfrac{0,4663}{2} = 0,2332 \ \text{mol}
        \end{equation*}

        \begin{equation*}
            m(\ch{Cu}) = 0,2332 \times 63,5 = \SI{14.8}{\gram}
        \end{equation*}

        \textbf{b) Volumen de \ch{Cl2} en condiciones normales}

        En el ánodo tiene lugar la oxidación \(\ch{2 Cl- -> Cl2 + 2 e-}\), por lo que:

        \begin{equation*}
            n(\ch{Cl2}) = \dfrac{n(e^-)}{2} = 0,2332 \ \text{mol}
        \end{equation*}

        \begin{equation*}
            V(\ch{Cl2}) = 0,2332 \times 22,4 = \SI{5.22}{\litre}
        \end{equation*}

    \end{solu}
\end{exercise}
\begin{exercise}
    Dos cubas electrolíticas que contienen disoluciones acuosas de nitrato de plata y nitrato de cobre (II), respectivamente están montadas en serie (pasa la misma intensidad por ambas). Si en 1 hora se depositan en la segunda cuba \masa[g]{54.5} de cobre, calcule:
    \begin{enumerate}[label=\letraminuscula]
        \item La intensidad de corriente que atraviesa las cubas.
        \item Los gramos de plata que se depositarán en la primera cuba tras dos horas de paso de la misma intensidad de corriente.
    \end{enumerate}
    \datos{} \(F = \num{96500}\ \dfrac{\text{C}}{\text{mol}}\); masa atómica Cu = \num{63.5}; masa atómica Ag = \num{107.9}
    \begin{solu}

        \textbf{a) Intensidad de corriente}

        En la segunda cuba (\ch{Cu(NO3)2}), la reducción es \(\ch{Cu^{2+} + 2 e- -> Cu}\). Los moles de \ch{Cu} depositados son:

        \begin{equation*}
            n(\ch{Cu}) = \dfrac{54,5}{63,5} = 0,8583 \ \text{mol} \ \Rightarrow \ n(e^-) = 2 \times 0,8583 = 1,7165 \ \text{mol}
        \end{equation*}

        \begin{equation*}
            Q = n(e^-) \cdot F = 1,7165 \times 96500 = \SI{165642}{\coulomb}
        \end{equation*}

        Como \(t = \SI{1}{\hour} = \SI{3600}{\s}\):

        \begin{equation*}
            I = \dfrac{Q}{t} = \dfrac{165642}{3600} = \SI{46.0}{\ampere}
        \end{equation*}

        \textbf{b) Gramos de plata depositados en la primera cuba tras 2 horas}

        Al estar las cubas en serie, la misma intensidad \(I = \SI{46.0}{\ampere}\) atraviesa ambas. Con \(t = \SI{2}{\hour} = \SI{7200}{\s}\):

        \begin{equation*}
            Q' = I \cdot t = 46,0 \times 7200 = \SI{331200}{\coulomb} \ \Rightarrow \ n(e^-) = \dfrac{331200}{96500} = 3,432 \ \text{mol}
        \end{equation*}

        Como en la primera cuba la reducción es \(\ch{Ag+ + e- -> Ag}\) (1 electrón por ion), los moles de \ch{Ag} depositados coinciden con los moles de electrones:

        \begin{equation*}
            n(\ch{Ag}) = 3,432 \ \text{mol} \ \Rightarrow \ m(\ch{Ag}) = 3,432 \times 107,9 = \SI{370.3}{\gram}
        \end{equation*}

    \end{solu}
    %* Libro Santillana página 307
\end{exercise}
\begin{exercise}
    Cuando se realiza la electrolisis en una disolución de una sal soluble de un metal divalente, pasando una corriente de \SI{3}{\ampere} durante cinco horas, se depositan \masa[g]{18.29}.
    \begin{enumerate}[label=\letraminuscula]
        \item Calcule la masa atómica del metal.
        \item Indique el nombre del electrodo donde se deposita el metal, escriba la reacción que tiene lugar en él y razone si se trata de una oxidación o de una reducción.
        \item Considerando el significado de la constante de Faraday, calcule expresando el resultado en culombios, la carga del ión divalente del metal.
    \end{enumerate}
    \datos{} \(F = \num{96500}\ \dfrac{\text{C}}{\text{mol}}\); \(N_A = \num{6.022e23}\).
    \begin{solu}

        \textbf{a) Masa atómica del metal}

        La carga total es:

        \begin{equation*}
            Q = I \cdot t = 3 \times (5 \times 3600) = \SI{54000}{\coulomb} \ \Rightarrow \ n(e^-) = \dfrac{54000}{96500} = 0,5596 \ \text{mol}
        \end{equation*}

        Como el metal es divalente, \(\ch{M^{2+} + 2 e- -> M}\):

        \begin{equation*}
            n(M) = \dfrac{n(e^-)}{2} = \dfrac{0,5596}{2} = 0,2798 \ \text{mol}
        \end{equation*}

        \begin{equation*}
            M_{at} = \dfrac{m}{n} = \dfrac{18,29}{0,2798} = \SI{65.4}{\gram\per\mole}
        \end{equation*}

        (Este valor corresponde, muy probablemente, al del cinc, \(M(\ch{Zn}) \approx \SI{65.4}{\gram\per\mole}\).)

        \textbf{b) Electrodo de deposición}

        El metal se deposita en el \textbf{cátodo}, donde tiene lugar la semirreacción:

        \begin{equation*}
            \ch{M^{2+} + 2 e- -> M}
        \end{equation*}

        Se trata de una \textbf{reducción}, ya que el catión metálico gana electrones para formar el metal neutro.

        \textbf{c) Carga del ion divalente}

        La constante de Faraday es la carga de 1 mol de electrones, \(F = N_A \cdot e\), por lo que la carga de un electrón es:

        \begin{equation*}
            e = \dfrac{F}{N_A} = \dfrac{96500}{6,022\times10^{23}} = \num{1.602e-19} \ \si{\coulomb}
        \end{equation*}

        Como el ion es divalente (carga \(2+\)), su carga es el doble de la del electrón (en valor absoluto):

        \begin{equation*}
            q(\ch{M^{2+}}) = 2 \times 1,602\times10^{-19} = \num{3.204e-19} \ \si{\coulomb}
        \end{equation*}

    \end{solu}
\end{exercise}

\subsection{Ejercicios PAU}
\begin{exerciseEVAU}
    Indique si las siguientes afirmaciones son verdaderas o falsas, justificando su
respuesta:
\begin{enumerate}[label=\letraminuscula]
    \item Cuanto mayor sea el \( E^0 \) de una sustancia, mayor será su carácter oxidante.
    \item La reacción de una pila en condiciones estándar transcurre sin necesidad de una fuente
de energía externa (proceso espontaneo), generando corriente eléctrica, si \( \Delta G > 0 \).
\end{enumerate}
\begin{solu}
    Indique si las siguientes afirmaciones son verdaderas o falsas, justificando su
respuesta:
letra
\end{solu}
\end{exerciseEVAU}
\begin{exerciseEVAU}
Sea la siguiente reacción redox:
\begin{equation*}
    \ch{Cu \sld{} + HNO3 -> Cu(NO3)2 + NO + H2O}
\end{equation*}
\begin{enumerate}[label=\letraminuscula]
    \item Escriba las semirreacciones de oxidación y reducción. Señale claramente cuál es el oxidante
y el reductor.
    \item  Ajuste las ecuaciones iónica y molecular por el método del ion-electrón.
    \item Calcule los gramos de Cu(s) necesarios para obtener 10 g de NO si el rendimiento de la
reacción es del 75\%.
\end{enumerate}
\begin{solu}

\end{solu}
\end{exerciseEVAU}

\begin{exerciseEVAU}
    Una forma de estimar la contaminación del agua es medir la Demanda Química de Oxígeno (DQO). Para
ello, el carbono (C) se transforma en dióxido de carbono (\ch{CO2}) al reaccionar en medio ácido con dicromato
de potasio (\ch{K2Cr2O7}), según la siguiente reacción:
    \begin{centeredequation}
        \ch{C + K2Cr2O7 + HCl -> CO2 + CrCl3 + KCl + H2O}
    \end{centeredequation}
    \begin{enumerate}[label=\letraminuscula]
        \item Escriba y ajuste las semirreacciones iónicas de oxidación y de reducción.
        \item Indique qué especie actúa como oxidante y cuál como reductor.
        \item Ajuste la reacción iónica global y la reacción molecular.
        \item Si \( E^0(\ch{Cr2O7^{2–}}/Cr^{3+}) \) = \qty{1.44}{\volt}  ¿qué valor debería tener \( E^00(\ch{CO2}/C) \) para que la reacción estuviera en equilibrio?
    \end{enumerate}
\begin{solu}

    \textbf{a) Semirreacciones}

    Oxidación (\ch{C^{0} -> C^{+4}}, pierde 4 e\textsuperscript{-}):

    \begin{equation*}
        \ch{C + 2 H2O -> CO2 + 4 H+ + 4 e-}
    \end{equation*}

    Reducción (\ch{Cr^{+6} -> Cr^{+3}}, gana \(3\times2=6\) e\textsuperscript{-} por \ch{Cr2O7^2-}):

    \begin{equation*}
        \ch{Cr2O7^2- + 14 H+ + 6 e- -> 2 Cr^{3+} + 7 H2O}
    \end{equation*}

    \textbf{b) Oxidante y reductor}

    El \textbf{oxidante} es el \ch{K2Cr2O7} (el \ch{Cr2O7^2-} se reduce), y el \textbf{reductor} es el \ch{C} (se oxida).

    \textbf{c) Reacción iónica y molecular}

    Igualando electrones (m.c.m. de 4 y 6 es 12): oxidación \(\times 3\), reducción \(\times 2\):

    \begin{equation*}
        \ch{3 C + 6 H2O -> 3 CO2 + 12 H+ + 12 e-}
    \end{equation*}
    \begin{equation*}
        \ch{2 Cr2O7^2- + 28 H+ + 12 e- -> 4 Cr^{3+} + 14 H2O}
    \end{equation*}

    Sumando y simplificando:

    \begin{equation*}
        \ch{3 C + 2 Cr2O7^2- + 16 H+ -> 3 CO2 + 4 Cr^{3+} + 8 H2O}
    \end{equation*}

    En forma molecular:

    \begin{equation*}
        \ch{3 C + 2 K2Cr2O7 + 16 HCl -> 3 CO2 + 4 CrCl3 + 4 KCl + 8 H2O}
    \end{equation*}

    \textbf{d) Potencial para el equilibrio}

    Una reacción redox está en equilibrio (\(\Delta E^0_{cell} = 0\), es decir, no transcurre en ningún sentido de forma espontánea) cuando los potenciales estándar de ambos pares son iguales:

    \begin{equation*}
        E^0(\ch{Cr2O7^{2-}}/\ch{Cr^{3+}}) = E^0(\ch{CO2}/\ch{C}) = \SI{1.44}{\volt}
    \end{equation*}

\end{solu}
\end{exerciseEVAU}

\begin{exerciseEVAU}
El dicromato de potasio \ch{K2Cr2O7} reacciona con cloruro de hierro (II), \ch{FeCl2}, en disolución de ácido
clorhídrico, \ch{HCl}, según la siguiente reacción:
\begin{centeredequation}
    \ch{K2Cr2O7 + FeCl2 + HCl -> CrCl3 + FeCl3 + KCl + H2O}
\end{centeredequation}
\begin{enumerate}[label=\letraminuscula]
    \item  Ajuste por el método del ion electrón las semirreacciones de oxidación y reducción. Indique
las especies oxidante y reductora. Ajuste la reacción iónica y la molecular.
    \item  Determine qué masa de \ch{K2Cr2O7} se necesitará para que reaccione completamente con 50 mL de disolución de \ch{FeCl2} 0,60 M
\end{enumerate}
\begin{solu}

    \textbf{a) Semirreacciones, oxidante/reductor y ecuaciones ajustadas}

    Reducción (\ch{Cr^{+6} -> Cr^{+3}}, gana 6 e\textsuperscript{-}):

    \begin{equation*}
        \ch{Cr2O7^2- + 14 H+ + 6 e- -> 2 Cr^{3+} + 7 H2O}
    \end{equation*}

    Oxidación (\ch{Fe^{+2} -> Fe^{+3}}, pierde 1 e\textsuperscript{-}), multiplicada por 6:

    \begin{equation*}
        \ch{6 Fe^{2+} -> 6 Fe^{3+} + 6 e-}
    \end{equation*}

    El \textbf{oxidante} es el \ch{K2Cr2O7} (\ch{Cr2O7^2-} se reduce) y el \textbf{reductor} es el \ch{FeCl2} (\ch{Fe^2+} se oxida).

    Reacción iónica global:

    \begin{equation*}
        \ch{Cr2O7^2- + 14 H+ + 6 Fe^{2+} -> 2 Cr^{3+} + 7 H2O + 6 Fe^{3+}}
    \end{equation*}

    Reacción molecular:

    \begin{equation*}
        \ch{K2Cr2O7 + 6 FeCl2 + 14 HCl -> 2 CrCl3 + 6 FeCl3 + 2 KCl + 7 H2O}
    \end{equation*}

    \textbf{b) Masa de \ch{K2Cr2O7} necesaria}

    \begin{equation*}
        n(\ch{FeCl2}) = 0,050 \times 0,60 = 0,03 \ \text{mol}
    \end{equation*}

    Por la estequiometría (1 mol \ch{K2Cr2O7} por cada 6 mol \ch{Fe^2+}):

    \begin{equation*}
        n(\ch{K2Cr2O7}) = \dfrac{0,03}{6} = 0,005 \ \text{mol}
    \end{equation*}

    Con \(M(\ch{K2Cr2O7}) = 2(39)+2(52)+7(16) = \SI{294}{\gram\per\mole}\):

    \begin{equation*}
        m(\ch{K2Cr2O7}) = 0,005 \times 294 = \SI{1.47}{\gram}
    \end{equation*}

\end{solu}
\end{exerciseEVAU}

\begin{exerciseEVAU}
    
El \ch{K2Cr2O7} reacciona con \ch{HI} en medio ácido sulfúrico para dar \ch{K2SO4}, \ch{Cr2(SO4)3}, \ch{I2} y \ch{H2O} según la siguiente reacción química:
    
    \ch{K2Cr2O7 + HI + H2SO4 -> K2SO4 + Cr2(SO4)3 + I2 + H2O}

    \begin{enumerate}[label=\letraminuscula]
        \item Escriba y ajuste las semirreacciones de oxidación y reducción, indicando la
especie que actúa como oxidante y la que actúa como reductor.
        \item  Escriba las reacciones iónica y molecular ajustadas. Utilice el método de ajuste de ión-electrón.
        \item Calcule la masa de \ch{I2} que se obtiene a partir de la reacción de 50 mL de dicromato potásico (\ch{K2Cr2O7}) 0,25 M, con exceso de ácido yodhídrico (\ch{HI}).
    \end{enumerate}

    \begin{solu}

        \textbf{a) Semirreacciones y especies oxidante/reductora}

        Reducción (\ch{Cr^{+6} -> Cr^{+3}}, gana 6 e\textsuperscript{-}):

        \begin{equation*}
            \ch{Cr2O7^2- + 14 H+ + 6 e- -> 2 Cr^{3+} + 7 H2O}
        \end{equation*}

        Oxidación (\ch{I^{-1} -> I^{0}}, pierde 1 e\textsuperscript{-} por átomo), multiplicada por 3:

        \begin{equation*}
            \ch{6 I- -> 3 I2 + 6 e-}
        \end{equation*}

        El \textbf{oxidante} es el \ch{K2Cr2O7} y el \textbf{reductor} es el \ch{HI}.

        \textbf{b) Reacciones iónica y molecular}

        \begin{equation*}
            \ch{Cr2O7^2- + 14 H+ + 6 I- -> 2 Cr^{3+} + 7 H2O + 3 I2}
        \end{equation*}

        \begin{equation*}
            \ch{K2Cr2O7 + 6 HI + 4 H2SO4 -> K2SO4 + Cr2(SO4)3 + 3 I2 + 7 H2O}
        \end{equation*}

        \textbf{c) Masa de \ch{I2} obtenida}

        \begin{equation*}
            n(\ch{K2Cr2O7}) = 0,050 \times 0,25 = 0,0125 \ \text{mol}
        \end{equation*}

        Por estequiometría (1 mol \ch{K2Cr2O7} produce 3 mol \ch{I2}):

        \begin{equation*}
            n(\ch{I2}) = 0,0125 \times 3 = 0,0375 \ \text{mol}
        \end{equation*}

        Con \(M(\ch{I2}) = \SI{254}{\gram\per\mole}\):

        \begin{equation*}
            m(\ch{I2}) = 0,0375 \times 254 = \SI{9.53}{\gram}
        \end{equation*}

    \end{solu}

\end{exerciseEVAU}

\begin{exerciseEVAU}

Dada la siguiente reacción entre el permanganato de potasio (\ch{KMnO4}) y el yoduro de potasio (\ch{KI}) en medio básico (\ch{KOH}) para dar manganato de potasio (\ch{K2MnO4}), iodato de potasio (\ch{KIO3}) y agua:
    \begin{enumerate}[label=\letraminuscula]
        \item Ajuste las semireacciones de oxidación y reducción e indique claramente la especie
oxidante y la especie reductora.
        \item Ajuste tanto la reacción iónica como la molecular según el método del ión-electrón.
        \item Calcule la cantidad de iodato de potasio (\ch{KIO3}) que se puede obtener a partir de 125 gramos de permanganato de potasio (\ch{KMnO4}) si el rendimiento de la operación es del 70 \%.
    \end{enumerate}

    \begin{solu}

        \textbf{a) Semirreacciones y especies oxidante/reductora}

        Reducción (\ch{Mn^{+7} -> Mn^{+6}}, gana 1 e\textsuperscript{-}):

        \begin{equation*}
            \ch{MnO4- + e- -> MnO4^2-}
        \end{equation*}

        Oxidación (\ch{I^{-1} -> I^{+5}}, pierde 6 e\textsuperscript{-}), en medio básico:

        \begin{equation*}
            \ch{I- + 6 OH- -> IO3- + 3 H2O + 6 e-}
        \end{equation*}

        El \textbf{oxidante} es el \ch{KMnO4} y el \textbf{reductor} es el \ch{KI}.

        \textbf{b) Reacciones iónica y molecular}

        Multiplicando la reducción por 6 para igualar los electrones:

        \begin{equation*}
            \ch{6 MnO4- + I- + 6 OH- -> 6 MnO4^2- + IO3- + 3 H2O}
        \end{equation*}

        \begin{equation*}
            \ch{6 KMnO4 + KI + 6 KOH -> 6 K2MnO4 + KIO3 + 3 H2O}
        \end{equation*}

        \textbf{c) Masa de \ch{KIO3} obtenida}

        Con \(M(\ch{KMnO4}) = 39+55+4(16) = \SI{158}{\gram\per\mole}\):

        \begin{equation*}
            n(\ch{KMnO4}) = \dfrac{125}{158} = 0,7911 \ \text{mol}
        \end{equation*}

        Por estequiometría (6 mol \ch{KMnO4} por cada 1 mol \ch{KIO3}), el rendimiento teórico es:

        \begin{equation*}
            n(\ch{KIO3})_{teórico} = \dfrac{0,7911}{6} = 0,1319 \ \text{mol}
        \end{equation*}

        Con un rendimiento del 70\,\%:

        \begin{equation*}
            n(\ch{KIO3})_{real} = 0,1319 \times 0,70 = 0,0923 \ \text{mol}
        \end{equation*}

        Con \(M(\ch{KIO3}) = 39+127+3(16) = \SI{214}{\gram\per\mole}\):

        \begin{equation*}
            m(\ch{KIO3}) = 0,0923 \times 214 = \SI{19.8}{\gram}
        \end{equation*}

    \end{solu}
\end{exerciseEVAU}

\begin{exerciseEVAU}

   En la siguiente reacción redox:

    \begin{reaction*}
        KMnO4 + KCl + H2SO4 -> MnSO4 + Cl2 + KHSO4 + H2O
    \end{reaction*}

    \begin{enumerate}[label=\letraminuscula]
        \item Escriba las semirreacciones de oxidación y reducción y señale claramente cuál es el oxidante
        y el reductor.
        \item Ajuste las ecuaciones iónica y molecular.
        \item Calcule cuántos gramos de KCl hacen falta para reaccionar completamente con \SI{5.5}{\g} de \ch{KMnO4}
        KMnO4
    \end{enumerate}

    Datos: Masas atómicas relativas: K = 39; Mn = 55; O = 16; Cl = 35,5

    \begin{solu}

        \textbf{a) Semirreacciones y especies oxidante/reductora}

        Reducción (\ch{Mn^{+7} -> Mn^{+2}}, gana 5 e\textsuperscript{-}):

        \begin{equation*}
            \ch{MnO4- + 8 H+ + 5 e- -> Mn^{2+} + 4 H2O}
        \end{equation*}

        Oxidación (\ch{Cl^{-1} -> Cl^{0}}, pierde 1 e\textsuperscript{-} por átomo):

        \begin{equation*}
            \ch{2 Cl- -> Cl2 + 2 e-}
        \end{equation*}

        El \textbf{oxidante} es el \ch{KMnO4} y el \textbf{reductor} es el \ch{KCl}.

        \textbf{b) Reacciones iónica y molecular}

        Igualando electrones (m.c.m. de 5 y 2 es 10): reducción \(\times 2\), oxidación \(\times 5\):

        \begin{equation*}
            \ch{2 MnO4- + 16 H+ + 10 Cl- -> 2 Mn^{2+} + 8 H2O + 5 Cl2}
        \end{equation*}

        En forma molecular, teniendo en cuenta que el exceso de sulfato queda como \ch{KHSO4}:

        \begin{equation*}
            \ch{2 KMnO4 + 10 KCl + 14 H2SO4 -> 2 MnSO4 + 5 Cl2 + 8 H2O + 12 KHSO4}
        \end{equation*}

        \textbf{c) Gramos de \ch{KCl} necesarios}

        Con \(M(\ch{KMnO4}) = 39+55+4(16) = \SI{158}{\gram\per\mole}\):

        \begin{equation*}
            n(\ch{KMnO4}) = \dfrac{5,5}{158} = 0,0348 \ \text{mol}
        \end{equation*}

        Por estequiometría (2 mol \ch{KMnO4} por cada 10 mol \ch{KCl}, es decir, relación 1:5):

        \begin{equation*}
            n(\ch{KCl}) = 0,0348 \times 5 = 0,1741 \ \text{mol}
        \end{equation*}

        Con \(M(\ch{KCl}) = 39+35,5 = \SI{74.5}{\gram\per\mole}\):

        \begin{equation*}
            m(\ch{KCl}) = 0,1741 \times 74,5 = \SI{13.0}{\gram}
        \end{equation*}

    \end{solu}
\end{exerciseEVAU}

\begin{exerciseEVAU}
     Se quiere construir una pila galvánica a partir de los pares \ch{Pb\pch[2]/ Pb} y \ch{Zn\pch[2]/Zn}.
    \begin{enumerate}[label=\letraminuscula]
        \item Indique cuál es el cátodo y cuál es el ánodo, cuál es el oxidante y cuál es el reductor.
        Escriba las reacciones que tienen lugar en cada uno de los electrodos
        \item Haga un esquema de la pila, indicando la dirección de circulación de los electrones
        y calcule el potencial estándar.
    \end{enumerate}
    \datos{}: Eº (\ch{Pb\pch[2]/ Pb}) = - 0,13 V; Eº (\ch{Zn\pch[2]/Zn}) = - 0,76 V
    \begin{solu}

        \textbf{a) Cátodo, ánodo, oxidante y reductor}

        Como \(E^0(\ch{Pb^{2+}}/\ch{Pb}) = -0,13\ \si{\volt} > E^0(\ch{Zn^{2+}}/\ch{Zn}) = -0,76\ \si{\volt}\), el par \ch{Pb^2+/Pb} tiene mayor tendencia a reducirse, por lo que actúa como \textbf{cátodo} (el \ch{Pb^2+} es el \textbf{oxidante}), y el par \ch{Zn^2+/Zn} actúa como \textbf{ánodo} (el \ch{Zn} es el \textbf{reductor}).

        \begin{equation*}
            \text{Cátodo (reducción):} \quad \ch{Pb^{2+} + 2 e- -> Pb}
        \end{equation*}
        \begin{equation*}
            \text{Ánodo (oxidación):} \quad \ch{Zn -> Zn^{2+} + 2 e-}
        \end{equation*}

        \textbf{b) Esquema y potencial estándar}

        La notación de la pila es \(\ch{Zn}(s) \mid \ch{Zn^{2+}}(aq) \parallel \ch{Pb^{2+}}(aq) \mid \ch{Pb}(s)\). Los electrones circulan por el circuito externo desde el ánodo (\ch{Zn}) hacia el cátodo (\ch{Pb}).

        \begin{equation*}
            E^0_{pila} = E^0_{cátodo} - E^0_{ánodo} = -0,13 - (-0,76) = \SI{0.63}{\volt}
        \end{equation*}

    \end{solu}
\end{exerciseEVAU}
\begin{exerciseEVAU}

    \begin{enumerate}[label=\letraminuscula]
        \item A la vista de los siguientes potenciales normales de reducción: Eº (\ch{H\pch[]} /\ch{H2}) = 0,00 V y Eº (\ch{Na\pch[]}/\ch{Na}) = -2,71 V, razone si se desprenderá hidrógeno (\ch{H2}) cuando se introduce una barra de sodio en una disolución
        de ácido clorhídrico.
        \item Se prepara una pila voltaica formada por electrodos estándar de Eº (\ch{Sn\pch[] / Sn}) = -0,1375 V y
        Eº (\ch{Pb\pch[]}/ \ch{Pb}) = -0,126 V. Escriba la semirreacción que ocurre en cada electrodo indicando cuál es el cátodo
        y cuál el ánodo, así como la reacción global ajustada.
        \item A la vista de los siguientes potenciales normales de reducción: Eº (\ch{Fe\pch[]}/\ch{Fe }) = -0,45 V y
        Eº (Cu2+ /Cu) = 0,34 V ¿Qué ocurrirá si se añaden limaduras de hierro (hierro en estado sólido) a una
        disolución de Cu2+? Justifique la respuesta indicando las semirreacciones que sucederán y la reacción global.
    \end{enumerate}
    \begin{solu}

        \textbf{a) Barra de sodio en ácido clorhídrico}

        Como \(E^0(\ch{Na+}/\ch{Na}) = \SI{-2.71}{\volt} \ll E^0(\ch{H+}/\ch{H2}) = \SI{0.00}{\volt}\), el sodio es un reductor mucho más fuerte que el hidrógeno, por lo que reducirá espontáneamente al \ch{H+}:

        \begin{equation*}
            E^0_{cell} = E^0(\ch{H+}/\ch{H2}) - E^0(\ch{Na+}/\ch{Na}) = 0,00-(-2,71) = \SI{2.71}{\volt} > 0
        \end{equation*}

        Por tanto, \textbf{sí se desprenderá hidrógeno} (reacción muy violenta): \(\ch{2 Na + 2 HCl -> 2 NaCl + H2}\).

        \textbf{b) Pila \ch{Sn}/\ch{Pb}}

        Como \(E^0(\ch{Pb^2+}/\ch{Pb}) = \SI{-0.126}{\volt} > E^0(\ch{Sn^2+}/\ch{Sn}) = \SI{-0.1375}{\volt}\), el plomo actúa como \textbf{cátodo} y el estaño como \textbf{ánodo}:

        \begin{equation*}
            \text{Cátodo:} \quad \ch{Pb^{2+} + 2 e- -> Pb} \qquad \text{Ánodo:} \quad \ch{Sn -> Sn^{2+} + 2 e-}
        \end{equation*}

        \begin{equation*}
            \ch{Pb^{2+} + Sn -> Pb + Sn^{2+}} \qquad E^0_{pila} = -0,126-(-0,1375) = \SI{0.0115}{\volt}
        \end{equation*}

        \textbf{c) Limaduras de hierro en disolución de \ch{Cu^2+}}

        Como \(E^0(\ch{Cu^2+}/\ch{Cu}) = \SI{0.34}{\volt} > E^0(\ch{Fe^2+}/\ch{Fe}) = \SI{-0.45}{\volt}\), el cobre se reduce (cátodo) y el hierro se oxida (ánodo):

        \begin{equation*}
            \text{Reducción:} \quad \ch{Cu^{2+} + 2 e- -> Cu} \qquad \text{Oxidación:} \quad \ch{Fe -> Fe^{2+} + 2 e-}
        \end{equation*}

        \begin{equation*}
            \ch{Fe + Cu^{2+} -> Fe^{2+} + Cu} \qquad E^0_{cell} = 0,34-(-0,45) = \SI{0.79}{\volt} > 0
        \end{equation*}

        Como \(E^0_{cell} > 0\), la reacción es espontánea: el hierro sólido se disuelve (se oxida a \ch{Fe^2+}) mientras se deposita cobre metálico sobre la limadura (reacción de cementación).

    \end{solu}
\end{exerciseEVAU}

\begin{exerciseEVAU}
    Sabiendo que los potenciales de reducción del cobre y de la plata en condiciones estándar son:
    Eo (\ch{Cu\pch[2]}/\ch{Cu}) = 0,35 V y Eo((\ch{Ag\pch[]}/\ch{Ag})) = 0,8 V:
    \begin{enumerate}[label=\letraminuscula]
        \item Calcule el potencial estándar de la pila que podría formarse con ellos.
        Indique razonadamente
        cuál sería el ánodo y cuál el cátodo.
        \item Atendiendo a la reacción global, conteste razonadamente si se disolverá una cuchara de
        cobre al introducirla en una disolución de \ch{AgNO3}.
        \item Escriba la notación de la pila formada en las condiciones del apartado a).
    \end{enumerate}
    \begin{solu}

        \textbf{a) Potencial estándar, ánodo y cátodo}

        Como \(E^0(\ch{Ag+}/\ch{Ag}) = \SI{0.8}{\volt} > E^0(\ch{Cu^2+}/\ch{Cu}) = \SI{0.35}{\volt}\), la plata actúa como \textbf{cátodo} (se reduce) y el cobre como \textbf{ánodo} (se oxida):

        \begin{equation*}
            E^0_{pila} = E^0_{cátodo} - E^0_{ánodo} = 0,8-0,35 = \SI{0.45}{\volt}
        \end{equation*}

        \textbf{b) ¿Se disuelve una cuchara de cobre en \ch{AgNO3}?}

        La reacción global es:

        \begin{equation*}
            \ch{Cu + 2 Ag+ -> Cu^{2+} + 2 Ag}
        \end{equation*}

        Como \(E^0_{pila} = \SI{0.45}{\volt} > 0\), la reacción es espontánea: \textbf{sí se disolverá} la cuchara de cobre (se oxida a \ch{Cu^2+}), mientras se deposita plata metálica sobre ella.

        \textbf{c) Notación de la pila}

        \begin{equation*}
            \ch{Cu}(s) \mid \ch{Cu^{2+}}(aq) \parallel \ch{Ag+}(aq) \mid \ch{Ag}(s)
        \end{equation*}

    \end{solu}
\end{exerciseEVAU}

\begin{exerciseEVAU}
    Cuando el yodo molecular, reacciona con el ácido nítrico se produce ácido yódico, dióxido de nitrógeno y agua.
    a) Escriba y ajuste, por el método del ion-electrón, las semirreacciones de oxidación y reducción que tienen lugar.
    Indique cuál es el oxidante y cuál el reductor.
    b) Escriba, ajustadas, la reacción iónica global y la reacción molecular global. c) Calcule el volumen de ácido nítrico del 65\% de riqueza en masa y densidad \densidad[gml]{1.5} que reacciona con \masa[g]{25.4}  de yodo molecular.
    Datos de masa atómica: I = 127, N = 14, O = 16, H = 1

    \begin{solu}

%        \begin{figure}[H]
%            \centering
%            \includegraphics[scale=.85]{images/imagesUD8ReaccionesRedox/UD8_Problemas/redox_EvAU_1.PNG}
%            %\caption{}
%            %\label{fig:}
%        \end{figure}

        \textbf{a) Semirreacciones}

        Oxidación (\ch{I^{0} -> I^{+5}}, pierde 5 e\textsuperscript{-} por átomo, 10 e\textsuperscript{-} por \ch{I2}):

        \begin{equation*}
            \ch{I2 + 6 H2O -> 2 IO3- + 12 H+ + 10 e-}
        \end{equation*}

        Reducción (\ch{N^{+5} -> N^{+4}}, gana 1 e\textsuperscript{-}):

        \begin{equation*}
            \ch{NO3- + 2 H+ + e- -> NO2 + H2O}
        \end{equation*}

        El \textbf{oxidante} es el \ch{HNO3} y el \textbf{reductor} es el \ch{I2}.

        \textbf{b) Reacciones iónica y molecular}

        Multiplicando la reducción por 10:

        \begin{equation*}
            \ch{I2 + 10 NO3- + 8 H+ -> 2 IO3- + 10 NO2 + 4 H2O}
        \end{equation*}

        \begin{equation*}
            \ch{I2 + 10 HNO3 -> 2 HIO3 + 10 NO2 + 4 H2O}
        \end{equation*}

        \textbf{c) Volumen de ácido nítrico}

        \begin{equation*}
            n(\ch{I2}) = \dfrac{25,4}{254} = 0,1 \ \text{mol}
        \end{equation*}

        Por estequiometría (1 mol \ch{I2} por cada 10 mol \ch{HNO3}):

        \begin{equation*}
            n(\ch{HNO3}) = 1 \ \text{mol} \ \Rightarrow \ m(\ch{HNO3}) = 1\times63 = \SI{63}{\gram}
        \end{equation*}

        Con una riqueza del 65\,\%:

        \begin{equation*}
            m_{disolución} = \dfrac{63}{0,65} = \SI{96.9}{\gram} \ \Rightarrow \ V = \dfrac{96,9}{1,5} = \SI{64.6}{\milli\litre}
        \end{equation*}

    \end{solu}
\end{exerciseEVAU}


\begin{exerciseEVAU}
    Para el siguiente proceso redox:

    \begin{equation*}
        \ch{K2Cr2O7 + HClO4 + HI -> KClO4 + Cr (ClO4)3 + I2 + H2O}
    \end{equation*}

    a) Escriba las semirreacciones de oxidación y reducción. Señale claramente cuál es el oxidante y el reductor. b) Ajuste las ecuaciones iónica y molecular por el método del ion-electrón. c) Calcule los gramos de \ch{K2Cr2O7} necesarios para obtener 60 g de \ch{I2} si el rendimiento de la reacción es del 50\%. Masas atómicas: K = 39,1; Cr = 52,0; H = 1,0; O = 16,0; Cl = 35,5; I =126,9

    \begin{solu}

%        \begin{figure}[H]
%            \centering
%            \includegraphics[scale=.8]{images/imagesUD8ReaccionesRedox/UD8_Problemas/redox_Evau_2.png}
%            %\caption{}
%            %\label{fig:}
%        \end{figure}

        \textbf{a) Semirreacciones}

        Reducción (\ch{Cr^{+6} -> Cr^{+3}}, gana 6 e\textsuperscript{-}):

        \begin{equation*}
            \ch{Cr2O7^2- + 14 H+ + 6 e- -> 2 Cr^{3+} + 7 H2O}
        \end{equation*}

        Oxidación (\ch{I^{-1} -> I^{0}}, pierde 1 e\textsuperscript{-} por átomo), multiplicada por 3:

        \begin{equation*}
            \ch{6 I- -> 3 I2 + 6 e-}
        \end{equation*}

        El \textbf{oxidante} es el \ch{K2Cr2O7} y el \textbf{reductor} es el \ch{HI}.

        \textbf{b) Reacciones iónica y molecular}

        \begin{equation*}
            \ch{Cr2O7^2- + 14 H+ + 6 I- -> 2 Cr^{3+} + 7 H2O + 3 I2}
        \end{equation*}

        \begin{equation*}
            \ch{K2Cr2O7 + 6 HI + 8 HClO4 -> 2 KClO4 + 2 Cr(ClO4)3 + 3 I2 + 7 H2O}
        \end{equation*}

        \textbf{c) Gramos de \ch{K2Cr2O7}}

        Con \(M(\ch{I2}) = 2(126,9) = \SI{253.8}{\gram\per\mole}\), y como el rendimiento es del 50\,\%, la cantidad de \ch{I2} teórica necesaria para obtener 60 g reales es:

        \begin{equation*}
            m(\ch{I2})_{teórica} = \dfrac{60}{0,50} = \SI{120}{\gram} \ \Rightarrow \ n(\ch{I2})_{teórica} = \dfrac{120}{253,8} = 0,4728 \ \text{mol}
        \end{equation*}

        Por estequiometría (1 mol \ch{K2Cr2O7} por cada 3 mol \ch{I2}):

        \begin{equation*}
            n(\ch{K2Cr2O7}) = \dfrac{0,4728}{3} = 0,1576 \ \text{mol}
        \end{equation*}

        Con \(M(\ch{K2Cr2O7}) = 2(39,1)+2(52,0)+7(16,0) = \SI{294.2}{\gram\per\mole}\):

        \begin{equation*}
            m(\ch{K2Cr2O7}) = 0,1576 \times 294,2 = \SI{46.4}{\gram}
        \end{equation*}

    \end{solu}
\end{exerciseEVAU}


\begin{exerciseEVAU}
    Para el siguiente proceso redox:
    \begin{equation*}
        \ch{KMnO4 + H2SO4 + KI -> MnSO4 + H2O + I2 + K2SO4}
    \end{equation*}
    a) Escriba las semirreacciones de oxidación y reducción. Señale claramente cuál es el oxidante y el reductor. b) Ajuste las ecuaciones iónica y molecular. c) Calcule los gramos de \ch{KMnO4} necesarios para obtener 30 g de \ch{I2}  si el rendimiento de la reacción es del 60\% Datos: Masas atómicas K = 39,1 Mn = 54,9 O = 16; I = 126,9

    \begin{solu}

%        \begin{figure}[H]
%            \centering
%            \includegraphics[scale=.8]{images/imagesUD8ReaccionesRedox/UD8_Problemas/redox_Evau_3.png}
%            %\caption{}
%            %\label{fig:}
%        \end{figure}

        \textbf{a) Semirreacciones}

        Reducción (\ch{Mn^{+7} -> Mn^{+2}}, gana 5 e\textsuperscript{-}):

        \begin{equation*}
            \ch{MnO4- + 8 H+ + 5 e- -> Mn^{2+} + 4 H2O}
        \end{equation*}

        Oxidación (\ch{I^{-1} -> I^{0}}, pierde 1 e\textsuperscript{-} por átomo), multiplicada por 5:

        \begin{equation*}
            \ch{10 I- -> 5 I2 + 10 e-}
        \end{equation*}

        El \textbf{oxidante} es el \ch{KMnO4} y el \textbf{reductor} es el \ch{KI}.

        \textbf{b) Reacciones iónica y molecular}

        Multiplicando la reducción por 2 (m.c.m. de 5 y 10):

        \begin{equation*}
            \ch{2 MnO4- + 16 H+ + 10 I- -> 2 Mn^{2+} + 8 H2O + 5 I2}
        \end{equation*}

        \begin{equation*}
            \ch{2 KMnO4 + 8 H2SO4 + 10 KI -> 2 MnSO4 + 8 H2O + 5 I2 + 6 K2SO4}
        \end{equation*}

        \textbf{c) Gramos de \ch{KMnO4}}

        Con \(M(\ch{I2}) = 2(126,9) = \SI{253.8}{\gram\per\mole}\) y rendimiento del 60\,\%, para obtener 30 g reales se necesita producir en teoría:

        \begin{equation*}
            n(\ch{I2})_{teórico} = \dfrac{30/0,60}{253,8} = \dfrac{50}{253,8} = 0,1970 \ \text{mol}
        \end{equation*}

        Por estequiometría (2 mol \ch{KMnO4} por cada 5 mol \ch{I2}):

        \begin{equation*}
            n(\ch{KMnO4}) = 0,1970 \times \dfrac{2}{5} = 0,0788 \ \text{mol}
        \end{equation*}

        Con \(M(\ch{KMnO4}) = 39,1+54,9+4(16) = \SI{158}{\gram\per\mole}\):

        \begin{equation*}
            m(\ch{KMnO4}) = 0,0788 \times 158 = \SI{12.4}{\gram}
        \end{equation*}

    \end{solu}
\end{exerciseEVAU}

\begin{exerciseEVAU}
    Para la siguiente reacción:
    \begin{equation*}
        \ch{KClO3 + SbCl3 + HCl → SbCl5 + KCl + H2O}
    \end{equation*}
    a) Nombrar todos los reactivos y productos, identificando razonadamente el oxidante y el reductor. b) Utilizando el método del ion electrón escribir y ajustar las semirreacciones de oxidación y reducción. c) Ajustar las ecuaciones iónica y molecular. d) Calcular cuántos gramos de \ch{KClO3} se necesitan para obtener 200 g de \ch{SbCl5}, si el rendimiento del proceso es del 70\%. Datos: Masas atómicas Sb = 121,8; Cl = 35,5; O = 16,0; K = 39,1.

    \solutionCM{a) clorato de potasio,  cloruro de antimonio (III), cloruro de hidrógeno -> cloruro de antimonio (V), cloruro de potasio, agua; reductor será el \ch{SbCl3} y el oxidante sería el \ch{ClO3-}; b) \negsubro{iónica}: \ch{ClO3- + 4 H+ + 3 Sb\pch[3] -> Cl- + 3 H2O + 3 Sb\pch[5]}. \textbf{molecular}: \ch{KClO3 + 3 SbCl3 + 6 HCl -> 3 SbCl5 + KCl + 3 H2O}; d) masa(\ch{KClO3} = \qty{38.53}{\g})  }
    \begin{solu}

%        \begin{figure}[H]
%            \centering
%            \includegraphics[scale=.8]{images/imagesUD8ReaccionesRedox/UD8_Problemas/redox_Evau_4.png}
%            %\caption{}
%            %\label{fig:}
%        \end{figure}

        \textbf{a) Nombres, oxidante y reductor}

        \ch{KClO3} (clorato de potasio), \ch{SbCl3} (cloruro de antimonio(III)) y \ch{HCl} (cloruro de hidrógeno) reaccionan para dar \ch{SbCl5} (cloruro de antimonio(V)), \ch{KCl} (cloruro de potasio) y \ch{H2O}. El \textbf{reductor} es el \ch{SbCl3} (el \ch{Sb^3+} se oxida a \ch{Sb^5+}) y el \textbf{oxidante} es el \ch{KClO3} (el \ch{ClO3-} se reduce a \ch{Cl-}).

        \textbf{b) Semirreacciones}

        Reducción (\ch{Cl^{+5} -> Cl^{-1}}, gana 6 e\textsuperscript{-}):

        \begin{equation*}
            \ch{ClO3- + 6 H+ + 6 e- -> Cl- + 3 H2O}
        \end{equation*}

        Oxidación (\ch{Sb^{+3} -> Sb^{+5}}, pierde 2 e\textsuperscript{-}), multiplicada por 3:

        \begin{equation*}
            \ch{3 Sb^{3+} -> 3 Sb^{5+} + 6 e-}
        \end{equation*}

        \textbf{c) Reacciones iónica y molecular}

        \begin{equation*}
            \ch{ClO3- + 6 H+ + 3 Sb^{3+} -> Cl- + 3 H2O + 3 Sb^{5+}}
        \end{equation*}

        \begin{equation*}
            \ch{KClO3 + 3 SbCl3 + 6 HCl -> 3 SbCl5 + KCl + 3 H2O}
        \end{equation*}

        \textbf{d) Gramos de \ch{KClO3} necesarios}

        Con \(M(\ch{SbCl5}) = 121,8+5(35,5) = \SI{299.3}{\gram\per\mole}\) y rendimiento del 70\,\%:

        \begin{equation*}
            n(\ch{SbCl5})_{teórico} = \dfrac{200/0,70}{299,3} = \dfrac{285,7}{299,3} = 0,9547 \ \text{mol}
        \end{equation*}

        Por estequiometría (1 mol \ch{KClO3} por cada 3 mol \ch{SbCl5}):

        \begin{equation*}
            n(\ch{KClO3}) = \dfrac{0,9547}{3} = 0,3182 \ \text{mol}
        \end{equation*}

        Con \(M(\ch{KClO3}) = 39,1+35,5+3(16,0) = \SI{122.6}{\gram\per\mole}\):

        \begin{equation*}
            m(\ch{KClO3}) = 0,3182 \times 122,6 = \SI{39.0}{\gram}
        \end{equation*}

    \end{solu}
\end{exerciseEVAU}

\begin{exerciseEVAU}
    Dada la reacción:
    \begin{equation*}
        \ch{Cu + HNO3 -> NO + Cu(NO3)2 + H2O}
    \end{equation*}
    \begin{enumerate}[label=\letraminuscula]
        \item Nombra los reactivos y los productos de la reacción. \solutionCM{Cobre + ácido nítrico -> monóxido de nitrico + nitrato de cobre (II) + agua}
        \item Utilizando el método del ion electrón escribe y ajusta las semiecuaciones de oxidación y reducción.
        \item ¿Qué volumen de \ch{HNO3} del 30\% (p/p) y densidad \densidad[gml]{1.18}1,18 g/mL se necesita para que reaccionen completamente \masa[g]{12.70} de \ch{Cu}? \solutionCM{\qty{94.32}{\mL}}
    \end{enumerate}
    \datos{}: Masas atómicas: Cu = 63,5; N = 14,0; H = 1,0; O = 16,0.

    \begin{solu}

%        \begin{figure}[H]
%            \centering
%            \includegraphics[scale=.8]{images/imagesUD8ReaccionesRedox/UD8_Problemas/redox_Evau_5.png}
%            %\caption{}
%            %\label{fig:}
%        \end{figure}

        \textbf{b) Semirreacciones}

        Reducción (\ch{N^{+5} -> N^{+2}}, gana 3 e\textsuperscript{-}):

        \begin{equation*}
            \ch{NO3- + 4 H+ + 3 e- -> NO + 2 H2O}
        \end{equation*}

        Oxidación (\ch{Cu^{0} -> Cu^{+2}}, pierde 2 e\textsuperscript{-}):

        \begin{equation*}
            \ch{Cu -> Cu^{2+} + 2 e-}
        \end{equation*}

        Igualando electrones (m.c.m. de 3 y 2 es 6): reducción \(\times 2\), oxidación \(\times 3\):

        \begin{equation*}
            \ch{3 Cu + 2 NO3- + 8 H+ -> 3 Cu^{2+} + 2 NO + 4 H2O}
        \end{equation*}

        En forma molecular:

        \begin{equation*}
            \ch{3 Cu + 8 HNO3 -> 3 Cu(NO3)2 + 2 NO + 4 H2O}
        \end{equation*}

        \textbf{c) Volumen de \ch{HNO3} necesario}

        \begin{equation*}
            n(\ch{Cu}) = \dfrac{12,70}{63,5} = 0,2 \ \text{mol}
        \end{equation*}

        Por estequiometría (3 mol \ch{Cu} por cada 8 mol \ch{HNO3}):

        \begin{equation*}
            n(\ch{HNO3}) = 0,2 \times \dfrac{8}{3} = 0,5333 \ \text{mol} \ \Rightarrow \ m(\ch{HNO3}) = 0,5333\times63 = \SI{33.6}{\gram}
        \end{equation*}

        Con una riqueza del 30\,\% (p/p) y \(\rho = \SI{1.18}{\gram\per\milli\litre}\):

        \begin{equation*}
            m_{disolución} = \dfrac{33,6}{0,30} = \SI{112.0}{\gram} \ \Rightarrow \ V = \dfrac{112,0}{1,18} = \SI{94.9}{\milli\litre}
        \end{equation*}

    \end{solu}
\end{exerciseEVAU}

\begin{exerciseEVAU}
    El \ch{KMnO4} reacciona con el \ch{KClO} en medio ácido \ch{H2SO4} dando \ch{MnSO4} y \ch{KClO3}, entre otros productos. a) Identifica y nombra el oxidante y el reductor. b) Ajusta la reacción por el método del ion electrón y escribe la ecuación iónica. c) Escribe la ecuación molecular ajustada. d) ¿Qué volumen de una disolución 0,05 M de KMnO4 será necesario para consumir \volumen[ml]{70} de disolución \concentracion[m]{0.02} de \ch{KClO}?

    \begin{solu}

%        \begin{figure}[H]
%            \centering
%            \includegraphics[scale=.8]{images/imagesUD8ReaccionesRedox/UD8_Problemas/redox_Evau_6.png}
%            %\caption{}
%            %\label{fig:}
%        \end{figure}

        \textbf{a) Oxidante y reductor}

        El \textbf{oxidante} es el \ch{KMnO4} (\ch{Mn^{+7} -> Mn^{+2}}, se reduce) y el \textbf{reductor} es el \ch{KClO} (\ch{Cl^{+1} -> Cl^{+5}}, se oxida a \ch{ClO3-}).

        \textbf{b) Semirreacciones y ecuación iónica}

        Reducción (gana 5 e\textsuperscript{-}):

        \begin{equation*}
            \ch{MnO4- + 8 H+ + 5 e- -> Mn^{2+} + 4 H2O}
        \end{equation*}

        Oxidación (pierde 4 e\textsuperscript{-}):

        \begin{equation*}
            \ch{ClO- + 2 H2O -> ClO3- + 4 H+ + 4 e-}
        \end{equation*}

        Igualando electrones (m.c.m. de 5 y 4 es 20): reducción \(\times 4\), oxidación \(\times 5\):

        \begin{equation*}
            \ch{4 MnO4- + 5 ClO- + 12 H+ -> 4 Mn^{2+} + 5 ClO3- + 6 H2O}
        \end{equation*}

        \textbf{c) Ecuación molecular}

        \begin{equation*}
            \ch{4 KMnO4 + 5 KClO + 6 H2SO4 -> 4 MnSO4 + 5 KClO3 + 2 K2SO4 + 6 H2O}
        \end{equation*}

        \textbf{d) Volumen de \ch{KMnO4} 0,05 M necesario}

        \begin{equation*}
            n(\ch{KClO}) = 0,070 \times 0,02 = \num{1.4e-3} \ \text{mol}
        \end{equation*}

        Por estequiometría (4 mol \ch{KMnO4} por cada 5 mol \ch{KClO}):

        \begin{equation*}
            n(\ch{KMnO4}) = 1,4\times10^{-3} \times \dfrac{4}{5} = \num{1.12e-3} \ \text{mol}
        \end{equation*}

        \begin{equation*}
            V(\ch{KMnO4}) = \dfrac{1,12\times10^{-3}}{0,05} = \num{2.24e-2} \ \si{\litre} = \SI{22.4}{\milli\litre}
        \end{equation*}

    \end{solu}
\end{exerciseEVAU}

\begin{exerciseEVAU}
    Dada la reacción:
    \begin{equation*}
        \ch{K2Cr2O7 + H2S + HCl CrCl3 + S + KCl + H2O}
    \end{equation*}
    a) Nombra los reactivos y los productos de la reacción. b) Utilizando el método del ion electrón escribe y ajusta las semiecuaciones de oxidación y reducción. c) Escribe las ecuaciones iónica y molecular ajustadas. d) Calcula cuánto azufre se produce si se consumen \masa[g]{51} de \ch{H2S} durante la reacción y esta transcurre con un rendimiento del 80\%. Datos: Masas atómicas S = 32,1; Cr = 52,0 K = 39,1; H = 1,0; O = 16,0

    \begin{solu}

%        \begin{figure}[H]
%            \centering
%            \includegraphics[scale=.8]{images/imagesUD8ReaccionesRedox/UD8_Problemas/redox_Evau_7.png}
%            %\caption{}
%            %\label{fig:}
%        \end{figure}

        \textbf{a) Nombres}

        \ch{K2Cr2O7} (dicromato de potasio), \ch{H2S} (sulfuro de hidrógeno) y \ch{HCl} (cloruro de hidrógeno) reaccionan para dar \ch{CrCl3} (cloruro de cromo(III)), \ch{S} (azufre) y \ch{H2O}, junto con \ch{KCl}.

        \textbf{b) Semirreacciones}

        Reducción (\ch{Cr^{+6} -> Cr^{+3}}, gana 6 e\textsuperscript{-}):

        \begin{equation*}
            \ch{Cr2O7^2- + 14 H+ + 6 e- -> 2 Cr^{3+} + 7 H2O}
        \end{equation*}

        Oxidación (\ch{S^{-2} -> S^{0}}, pierde 2 e\textsuperscript{-}), multiplicada por 3:

        \begin{equation*}
            \ch{3 H2S -> 3 S + 6 H+ + 6 e-}
        \end{equation*}

        \textbf{c) Reacciones iónica y molecular}

        \begin{equation*}
            \ch{Cr2O7^2- + 3 H2S + 8 H+ -> 2 Cr^{3+} + 3 S + 7 H2O}
        \end{equation*}

        \begin{equation*}
            \ch{K2Cr2O7 + 3 H2S + 8 HCl -> 2 CrCl3 + 3 S + 2 KCl + 7 H2O}
        \end{equation*}

        \textbf{d) Gramos de azufre producidos}

        Con \(M(\ch{H2S}) = 2(1,0)+32,1 = \SI{34.1}{\gram\per\mole}\):

        \begin{equation*}
            n(\ch{H2S}) = \dfrac{51}{34,1} = 1,4956 \ \text{mol}
        \end{equation*}

        Por estequiometría (3 mol \ch{H2S} por cada 3 mol \ch{S}, relación 1:1), y con un rendimiento del 80\,\%:

        \begin{equation*}
            n(\ch{S})_{real} = 1,4956 \times 0,80 = 1,1965 \ \text{mol}
        \end{equation*}

        \begin{equation*}
            m(\ch{S}) = 1,1965 \times 32,1 = \SI{38.4}{\gram}
        \end{equation*}

    \end{solu}
\end{exerciseEVAU}


\begin{exerciseEVAU}
    Se puede producir gas cloro haciendo uso de la siguiente reacción:
    \begin{equation*}
        \ch{K2Cr2O7 + HCl -> KCl + CrCl3 + Cl2 + H2O}
    \end{equation*}

    a) Escribe y ajusta las semirreacciones de oxidación y reducción utilizando el método del ion-electrón. Indica el nombre del oxidante y del reductor. b) Ajusta la ecuación molecular. c) Calcula los moles de \ch{Cl2} que se producirán si se consumen totalmente \masa[g]{18.25} de \ch{HCl}. DATOS: Masas atómicas: Cl = 35,5 u; H = 1,0 u.

    \begin{solu}

%        \begin{figure}[H]
%            \centering
%            \includegraphics[scale=.8]{images/imagesUD8ReaccionesRedox/UD8_Problemas/redox_Evau_8.png}
%            %\caption{}
%            %\label{fig:}
%        \end{figure}

        \textbf{a) Semirreacciones, oxidante y reductor}

        Reducción (\ch{Cr^{+6} -> Cr^{+3}}, gana 6 e\textsuperscript{-}):

        \begin{equation*}
            \ch{Cr2O7^2- + 14 H+ + 6 e- -> 2 Cr^{3+} + 7 H2O}
        \end{equation*}

        Oxidación (\ch{Cl^{-1} -> Cl^{0}}, pierde 1 e\textsuperscript{-} por átomo), multiplicada por 3:

        \begin{equation*}
            \ch{6 Cl- -> 3 Cl2 + 6 e-}
        \end{equation*}

        El \textbf{oxidante} es el \ch{K2Cr2O7} y el \textbf{reductor} es el \ch{HCl}.

        \textbf{b) Ecuación molecular}

        \begin{equation*}
            \ch{K2Cr2O7 + 14 HCl -> 2 KCl + 2 CrCl3 + 3 Cl2 + 7 H2O}
        \end{equation*}

        \textbf{c) Moles de \ch{Cl2} producidos}

        Con \(M(\ch{HCl}) = 35,5+1,0 = \SI{36.5}{\gram\per\mole}\):

        \begin{equation*}
            n(\ch{HCl}) = \dfrac{18,25}{36,5} = 0,5 \ \text{mol}
        \end{equation*}

        Por estequiometría (14 mol \ch{HCl} por cada 3 mol \ch{Cl2}):

        \begin{equation*}
            n(\ch{Cl2}) = 0,5 \times \dfrac{3}{14} = \num{0.107} \ \text{mol}
        \end{equation*}

    \end{solu}
\end{exerciseEVAU}

\begin{exerciseEVAU}
    Para la reacción de cobre con ácido nítrico que produce nitrato de cobre (II), monóxido de nitrógeno y agua. a) Ajusta la ecuación iónica y molecular por el método del ion-electrón. b) Señala el oxidante y el reductor. c) ¿Cuántos gramos de ácido nítrico son necesarios para obtener 5 L de óxido de nitrógeno medidos en condiciones normales?

    \datos{}: \ctegases; masas atómicas: H=1, N= 14, O= 16.

    \begin{solu}

%        \begin{figure}[H]
%            \centering
%            \includegraphics[scale=.8]{images/imagesUD8ReaccionesRedox/UD8_Problemas/redox_Evau_9.png}
%            %\caption{}
%            %\label{fig:}
%        \end{figure}

        \textbf{a) Semirreacciones, oxidante y reductor}

        Reducción (\ch{N^{+5} -> N^{+2}}, gana 3 e\textsuperscript{-}):

        \begin{equation*}
            \ch{NO3- + 4 H+ + 3 e- -> NO + 2 H2O}
        \end{equation*}

        Oxidación (\ch{Cu^{0} -> Cu^{+2}}, pierde 2 e\textsuperscript{-}):

        \begin{equation*}
            \ch{Cu -> Cu^{2+} + 2 e-}
        \end{equation*}

        El \textbf{oxidante} es el \ch{HNO3} y el \textbf{reductor} es el \ch{Cu}.

        \textbf{b) Ecuaciones iónica y molecular}

        Igualando electrones (m.c.m. de 3 y 2 es 6): reducción \(\times 2\), oxidación \(\times 3\):

        \begin{equation*}
            \ch{3 Cu + 2 NO3- + 8 H+ -> 3 Cu^{2+} + 2 NO + 4 H2O}
        \end{equation*}

        \begin{equation*}
            \ch{3 Cu + 8 HNO3 -> 3 Cu(NO3)2 + 2 NO + 4 H2O}
        \end{equation*}

        \textbf{c) Gramos de \ch{HNO3} necesarios}

        Los moles de \ch{NO} en 5 L a condiciones normales son:

        \begin{equation*}
            n(\ch{NO}) = \dfrac{5}{22,4} = 0,2232 \ \text{mol}
        \end{equation*}

        Por estequiometría (8 mol \ch{HNO3} por cada 2 mol \ch{NO}):

        \begin{equation*}
            n(\ch{HNO3}) = 0,2232 \times \dfrac{8}{2} = 0,8929 \ \text{mol}
        \end{equation*}

        Con \(M(\ch{HNO3}) = 1+14+3(16) = \SI{63}{\gram\per\mole}\):

        \begin{equation*}
            m(\ch{HNO3}) = 0,8929 \times 63 = \SI{56.3}{\gram}
        \end{equation*}

    \end{solu}
\end{exerciseEVAU}

\begin{exerciseEVAU}
    El dicromato de potasio es un oxidante fuerte que se utiliza en algunos preparados para proteger la madera. Este compuesto reacciona con el yoduro de potasio y el ácido sulfúrico, obteniéndose como productos sulfato de cromo (III), yodo molecular, sulfato de potasio y agua. a) Ajusta la ecuación iónica y molecular por el método del ión-electrón. b) Escribe las parejas de oxidante-reductor y oxidado-reducido. c) Sabiendo que al reaccionar \volumen[ml]{157}  de una disolución de dicromato de potasio, con suficiente yoduro de potasio y ácido sulfúrico, se obtienen \masa[g]{7.62} de yodo molecular, calcula la concentración de dicha disolución.
    \datos{}: masas atómicas: I = 25.

    \begin{solu}

%        \begin{figure}[H]
%            \centering
%            \includegraphics[scale=.8]{images/imagesUD8ReaccionesRedox/UD8_Problemas/redox_Evau_10.png}
%            %\caption{}
%            %\label{fig:}
%        \end{figure}

        \textbf{a) Ecuaciones iónica y molecular}

        Reducción (\ch{Cr^{+6} -> Cr^{+3}}, gana 6 e\textsuperscript{-}):

        \begin{equation*}
            \ch{Cr2O7^2- + 14 H+ + 6 e- -> 2 Cr^{3+} + 7 H2O}
        \end{equation*}

        Oxidación (\ch{I^{-1} -> I^{0}}, pierde 1 e\textsuperscript{-} por átomo), multiplicada por 3:

        \begin{equation*}
            \ch{6 I- -> 3 I2 + 6 e-}
        \end{equation*}

        \begin{equation*}
            \ch{Cr2O7^2- + 14 H+ + 6 I- -> 2 Cr^{3+} + 7 H2O + 3 I2}
        \end{equation*}

        \begin{equation*}
            \ch{K2Cr2O7 + 6 KI + 7 H2SO4 -> Cr2(SO4)3 + 3 I2 + 4 K2SO4 + 7 H2O}
        \end{equation*}

        \textbf{b) Pares oxidante-reductor y oxidado-reducido}

        Par oxidante/reducido: \(\ch{Cr2O7^2-}\) (oxidante) / \(\ch{Cr^{3+}}\) (reducido). Par reductor/oxidado: \(\ch{I-}\) (reductor) / \(\ch{I2}\) (oxidado).

        \textbf{c) Concentración de la disolución de dicromato}

        (Se emplea la masa atómica habitual del yodo, \(M(\ch{I}) \approx \SI{127}{\gram\per\mole}\), ya que el dato «I = 25» del enunciado es, con toda probabilidad, una errata.)

        \begin{equation*}
            n(\ch{I2}) = \dfrac{7,62}{2\times127} = 0,0300 \ \text{mol}
        \end{equation*}

        Por estequiometría (1 mol \ch{K2Cr2O7} por cada 3 mol \ch{I2}):

        \begin{equation*}
            n(\ch{K2Cr2O7}) = \dfrac{0,0300}{3} = 0,0100 \ \text{mol}
        \end{equation*}

        \begin{equation*}
            [\ch{K2Cr2O7}] = \dfrac{0,0100}{0,157} = \SI{0.0637}{\Molar}
        \end{equation*}

    \end{solu}
\end{exerciseEVAU}

\subsection{Pilas}

\begin{exerciseEVAU}
    Teniendo en cuenta los potenciales de reducción estándar de la plata y del níquel: \( E\text{º}(\ch{Ag+}\ |\ \ch{Ag}) = \SI{+80}{\volt} \) y \( E\text{º}(\ch{Ni^{2+}}\ |\  \ch{Ni}) = \SI{-0.25}{\volt} \): a) ¿Cuál es la fuerza electromotriz, en condiciones estándar, de la pila que se podría construir? b) Escriba la notación de esa pila y la reacción global que tienen lugar. c) Indique cuál es el cátodo, el ánodo, las reacciones que tienen lugar en cada uno de ellos.

    \begin{solu}

%        \begin{figure}[H]
%            \centering
%            \includegraphics[scale=.8]{images/imagesUD8ReaccionesRedox/UD8_Problemas/redox_Evau_11.png}
%            %\caption{}
%            %\label{fig:}
%        \end{figure}

        Como \(E^0(\ch{Ag+}/\ch{Ag}) = \SI{0.80}{\volt} > E^0(\ch{Ni^2+}/\ch{Ni}) = \SI{-0.25}{\volt}\), la plata actúa como cátodo y el níquel como ánodo.

        \textbf{a) Fuerza electromotriz}

        \begin{equation*}
            E^0_{pila} = E^0_{cátodo} - E^0_{ánodo} = 0,80-(-0,25) = \SI{1.05}{\volt}
        \end{equation*}

        \textbf{b) Notación y reacción global}

        \begin{equation*}
            \ch{Ni}(s) \mid \ch{Ni^{2+}}(aq) \parallel \ch{Ag+}(aq) \mid \ch{Ag}(s)
        \end{equation*}

        \begin{equation*}
            \ch{Ni + 2 Ag+ -> Ni^{2+} + 2 Ag}
        \end{equation*}

        \textbf{c) Cátodo, ánodo y semirreacciones}

        \begin{equation*}
            \text{Cátodo (reducción):} \quad \ch{2 Ag+ + 2 e- -> 2 Ag}
        \end{equation*}
        \begin{equation*}
            \text{Ánodo (oxidación):} \quad \ch{Ni -> Ni^{2+} + 2 e-}
        \end{equation*}

    \end{solu}
\end{exerciseEVAU}


\begin{exerciseEVAU}
    Sabiendo que los potenciales de reducción del hierro y del cobalto en condiciones estándar son \( E\text{º}(\ch{Fe\pch[2]}\ | \ \ch{Fe}) = \SI{-0.44}{\volt} \)  y \( E\text{º}(\ch{Co\pch[2]}\ | \ \ch{Co}) = \SI{-0.28}{\volt} \): a) Calcule el potencial estándar de la pila que podría formarse con ellos. Indique razonadamente cuál sería el ánodo y cuál el cátodo. b) Escriba las reacciones que tendrían lugar en el ánodo y en el cátodo, así como la reacción global de la pila. c) Escriba la notación de la pila.

    \begin{solu}

%        \begin{figure}[H]
%            \centering
%            \includegraphics[scale=.8]{images/imagesUD8ReaccionesRedox/UD8_Problemas/redox_Evau_12.png}
%            %\caption{}
%            %\label{fig:}
%        \end{figure}

        \textbf{a) Potencial estándar, ánodo y cátodo}

        Como \(E^0(\ch{Co^2+}/\ch{Co}) = \SI{-0.28}{\volt} > E^0(\ch{Fe^2+}/\ch{Fe}) = \SI{-0.44}{\volt}\), el cobalto actúa como \textbf{cátodo} y el hierro como \textbf{ánodo}:

        \begin{equation*}
            E^0_{pila} = -0,28-(-0,44) = \SI{0.16}{\volt}
        \end{equation*}

        \textbf{b) Reacciones en ánodo, cátodo y global}

        \begin{equation*}
            \text{Ánodo:} \quad \ch{Fe -> Fe^{2+} + 2 e-} \qquad \text{Cátodo:} \quad \ch{Co^{2+} + 2 e- -> Co}
        \end{equation*}

        \begin{equation*}
            \ch{Fe + Co^{2+} -> Fe^{2+} + Co}
        \end{equation*}

        \textbf{c) Notación de la pila}

        \begin{equation*}
            \ch{Fe}(s) \mid \ch{Fe^{2+}}(aq) \parallel \ch{Co^{2+}}(aq) \mid \ch{Co}(s)
        \end{equation*}

    \end{solu}
\end{exerciseEVAU}

\begin{exerciseEVAU}
    Sabiendo que los potenciales de reducción del cobre y de la plata en condiciones estándar son \( E\text{º}(\ch{Cu+}\ | \ \ch{Cu}) = \SI{+0.52}{\volt} \)  y \( E\text{º}(\ch{Ag+}\ | \ \ch{Ag}) = \SI{+0.8}{\volt} \): a) Indique razonadamente cuál sería el ánodo y cuál el cátodo. Calcule el potencial estándar de la pila que podría formarse con ellos. b) Escriba las reacciones que tendrían lugar en el ánodo y en el cátodo, así como la reacción global de la pila. c) Escriba la notación de la pila.

    \begin{solu}

%        \begin{figure}[H]
%            \centering
%            \includegraphics[scale=.8]{images/imagesUD8ReaccionesRedox/UD8_Problemas/redox_Evau_13.png}
%            %\caption{}
%            %\label{fig:}
%        \end{figure}

        \textbf{a) Ánodo, cátodo y potencial estándar}

        Como \(E^0(\ch{Ag+}/\ch{Ag}) = \SI{0.8}{\volt} > E^0(\ch{Cu+}/\ch{Cu}) = \SI{0.52}{\volt}\), la plata actúa como \textbf{cátodo} y el cobre como \textbf{ánodo}:

        \begin{equation*}
            E^0_{pila} = 0,8-0,52 = \SI{0.28}{\volt}
        \end{equation*}

        \textbf{b) Reacciones en ánodo, cátodo y global}

        \begin{equation*}
            \text{Ánodo:} \quad \ch{Cu -> Cu+ + e-} \qquad \text{Cátodo:} \quad \ch{Ag+ + e- -> Ag}
        \end{equation*}

        \begin{equation*}
            \ch{Cu + Ag+ -> Cu+ + Ag}
        \end{equation*}

        \textbf{c) Notación de la pila}

        \begin{equation*}
            \ch{Cu}(s) \mid \ch{Cu+}(aq) \parallel \ch{Ag+}(aq) \mid \ch{Ag}(s)
        \end{equation*}

    \end{solu}
\end{exerciseEVAU}

\begin{exerciseEVAU}
    Según la reacción:
    \begin{equation*}
        \ch{HNO3 + Mg -> NO2 + Mg (NO3)2 + H2O}
    \end{equation*}
    a) Escriba las semirreacciones de oxidación y de reducción. Indique cuál es la especie oxidante y cuál la reductora. b) Ajuste las reacciones iónica y molecular por el método del ion-electrón. c) Calcule el potencial de la pila. Datos: \( E\text{º}(\ch{Mg\pch[2]}\ | \ \ch{Mg}) = \SI{-2.37}{\volt} \); \( E\text{º}(\ch{NO3-}\ | \ \ch{NO2}) = \SI{+0.78}{\volt} \)

    \begin{solu}

%        \begin{figure}[H]
%            \centering
%            \includegraphics[scale=.8]{images/imagesUD8ReaccionesRedox/UD8_Problemas/redox_Evau_14.png}
%            %\caption{}
%            %\label{fig:}
%        \end{figure}

        \textbf{a) Semirreacciones, oxidante y reductor}

        Reducción (\ch{N^{+5} -> N^{+4}}, gana 1 e\textsuperscript{-}):

        \begin{equation*}
            \ch{NO3- + 2 H+ + e- -> NO2 + H2O}
        \end{equation*}

        Oxidación (\ch{Mg^{0} -> Mg^{+2}}, pierde 2 e\textsuperscript{-}):

        \begin{equation*}
            \ch{Mg -> Mg^{2+} + 2 e-}
        \end{equation*}

        El \textbf{oxidante} es el \ch{HNO3} y el \textbf{reductor} es el \ch{Mg}.

        \textbf{b) Reacciones iónica y molecular}

        Multiplicando la reducción por 2:

        \begin{equation*}
            \ch{Mg + 2 NO3- + 4 H+ -> Mg^{2+} + 2 NO2 + 2 H2O}
        \end{equation*}

        \begin{equation*}
            \ch{Mg + 4 HNO3 -> Mg(NO3)2 + 2 NO2 + 2 H2O}
        \end{equation*}

        \textbf{c) Potencial de la pila}

        Como \(E^0(\ch{NO3-}/\ch{NO2}) = \SI{0.78}{\volt} > E^0(\ch{Mg^2+}/\ch{Mg}) = \SI{-2.37}{\volt}\), el par \ch{NO3-}/\ch{NO2} actúa como cátodo:

        \begin{equation*}
            E^0_{pila} = 0,78-(-2,37) = \SI{3.15}{\volt}
        \end{equation*}

    \end{solu}
\end{exerciseEVAU}

\begin{exerciseEVAU}
    Teniendo en cuenta los potenciales de reducción estándar de los pares \( E\text{º}(\ch{Ag+}\ | \ \ch{Ag}) = \SI{+0.8}{\volt} \)  y \( E\text{º}(\ch{Ni\pch[2]}\ | \ \ch{Ni}) = \SI{-0.25}{\volt} \): a) ¿Cuál es la fuerza electromotriz, en condiciones estándar, de la pila que se podría construir? b) Escribir la notación de la pila y las reacciones anódica, catódica y global que tendrían lugar.

    \begin{solu}

%        \begin{figure}[H]
%            \centering
%            \includegraphics[scale=.8]{images/imagesUD8ReaccionesRedox/UD8_Problemas/redox_Evau_15.png}
%            %\caption{}
%            %\label{fig:}
%        \end{figure}

        Como \(E^0(\ch{Ag+}/\ch{Ag}) = \SI{0.8}{\volt} > E^0(\ch{Ni^2+}/\ch{Ni}) = \SI{-0.25}{\volt}\), la plata actúa como cátodo y el níquel como ánodo.

        \textbf{a) Fuerza electromotriz}

        \begin{equation*}
            E^0_{pila} = 0,8-(-0,25) = \SI{1.05}{\volt}
        \end{equation*}

        \textbf{b) Notación y reacciones}

        \begin{equation*}
            \ch{Ni}(s) \mid \ch{Ni^{2+}}(aq) \parallel \ch{Ag+}(aq) \mid \ch{Ag}(s)
        \end{equation*}

        \begin{equation*}
            \text{Ánodo:} \quad \ch{Ni -> Ni^{2+} + 2 e-} \qquad \text{Cátodo:} \quad \ch{2 Ag+ + 2 e- -> 2 Ag}
        \end{equation*}

        \begin{equation*}
            \text{Global:} \quad \ch{Ni + 2 Ag+ -> Ni^{2+} + 2 Ag}
        \end{equation*}

    \end{solu}
\end{exerciseEVAU}

\begin{exerciseEVAU}
    Se dispone en el laboratorio de dos barras metálicas, una de \ch{Cd} y otra de \ch{Al}, y de una disolución de \ch{Zn^{2+}} de concentración \SI{1}{\m}. a) Justifica cuál de las dos barras deberá introducirse en la disolución para obtener \ch{Zn} metálico. b) Escribe las semirreacciones de oxidación y reducción y la reacción global que tendrá lugar. Datos. \( E\text{º}(\ch{Cd\pch[2]}\ | \ \ch{Cd}) = \SI{-0.40}{\volt} \); \( E\text{º}(\ch{Zn\pch[2]}\ | \ \ch{Zn}) = \SI{-0.77}{\volt} \); \( E\text{º}(\ch{Al\pch[3]}\ | \ \ch{Al}) = \SI{-1.68}{\volt} \).

    \begin{solu}

%        \begin{figure}[H]
%            \centering
%            \includegraphics[scale=.8]{images/imagesUD8ReaccionesRedox/UD8_Problemas/redox_Evau_16.png}
%            %\caption{}
%            %\label{fig:}
%        \end{figure}

        \textbf{a) Barra a introducir}

        Para obtener \ch{Zn} metálico, la especie \(\ch{Zn^2+}\) debe reducirse, lo que requiere un reductor con potencial de reducción más negativo que \(E^0(\ch{Zn^2+}/\ch{Zn}) = \SI{-0.77}{\volt}\).

        \begin{itemize}
            \item \(E^0(\ch{Cd^2+}/\ch{Cd}) = \SI{-0.40}{\volt}\): menos negativo que el del \ch{Zn}, por lo que el \ch{Cd} \textbf{no} puede reducir al \ch{Zn^2+}.
            \item \(E^0(\ch{Al^3+}/\ch{Al}) = \SI{-1.68}{\volt}\): más negativo que el del \ch{Zn}, por lo que el \ch{Al} \textbf{sí} puede reducir al \ch{Zn^2+}.
        \end{itemize}

        Por tanto, debe introducirse la \textbf{barra de aluminio}.

        \textbf{b) Semirreacciones y reacción global}

        Oxidación (\(\times 2\)): \(\ch{Al -> Al^{3+} + 3 e-}\). Reducción (\(\times 3\)): \(\ch{Zn^{2+} + 2 e- -> Zn}\).

        \begin{equation*}
            \ch{2 Al + 3 Zn^{2+} -> 2 Al^{3+} + 3 Zn}
        \end{equation*}

        \begin{equation*}
            E^0_{pila} = E^0(\ch{Zn^2+}/\ch{Zn}) - E^0(\ch{Al^3+}/\ch{Al}) = -0,77-(-1,68) = \SI{0.91}{\volt} > 0
        \end{equation*}

        Al ser \(E^0_{pila} > 0\), se confirma que la reacción es espontánea con la barra de aluminio.

    \end{solu}
\end{exerciseEVAU}

\begin{exerciseEVAU}
    Se dispone de dos barras metálicas, una de plata y otra de cadmio, y de \volumen[ml]{100} de sendas disoluciones \SI{1}{\m}  de \ch{AgNO3} y \ch{Cd(NO3)2}. Justifica qué barra metálica habría que introducir en qué disolución para que se produzca una reacción espontánea. Datos: \( E\text{º}(\ch{Ag+}\ | \ \ch{Ag}) = \SI{+0.80}{\volt} \); \( E\text{º}(\ch{Cd\pch[2]}\ | \ \ch{Cd}) = \SI{-0.40}{\volt} \).

    \begin{solu}

%        \begin{figure}[H]
%            \centering
%            \includegraphics[scale=.55]{images/imagesUD8ReaccionesRedox/UD8_Problemas/redox_Evau_17.png}
%            %\caption{}
%            %\label{fig:}
%        \end{figure}

        Como \(E^0(\ch{Ag+}/\ch{Ag}) = \SI{0.80}{\volt} > E^0(\ch{Cd^2+}/\ch{Cd}) = \SI{-0.40}{\volt}\), la plata (\ch{Ag+}) es más fácil de reducir que el cadmio (\ch{Cd^2+}), y el \ch{Cd} es un reductor más fuerte que la \ch{Ag}.

        Para que se produzca una reacción espontánea, debe introducirse la \textbf{barra de cadmio en la disolución de \ch{AgNO3}}: el \ch{Cd} se oxida (ánodo) mientras el \ch{Ag+} se reduce (cátodo), depositándose plata metálica:

        \begin{equation*}
            \ch{Cd + 2 Ag+ -> Cd^{2+} + 2 Ag} \qquad E^0_{pila} = 0,80-(-0,40) = \SI{1.20}{\volt} > 0
        \end{equation*}

        En cambio, si se introdujera la barra de plata en la disolución de \ch{Cd(NO3)2}, el proceso inverso tendría \(E^0_{pila} = \SI{-1.20}{\volt} < 0\), por lo que \textbf{no sería espontáneo}.

    \end{solu}
\end{exerciseEVAU}

\begin{exerciseEVAU}
    Indica razonadamente si los siguientes procesos de oxidación-reducción pueden tener lugar de forma espontánea:
    \begin{enumerate}[label=\letraminuscula]
        \item \ch{Fe\pch[2] + 1/2 Cl2 -> Fe\pch[3] + Cl-}
        \item \ch{Fe\pch[2] + 1/2 Br2 -> Fe\pch[3] + Br-}
    \end{enumerate}
    \datos{}: \( E\text{º}(\ch{Fe\pch[2]}\ | \ \ch{Fe\pch[2]}) = \SI{0.77}{\volt} \); \( E\text{º}(\ch{Cl2}\ | \ \ch{Cl-}) = \SI{1.36}{\volt} \); \( E\text{º}(\ch{Br2}\ | \ \ch{Br-}) = \SI{0.54}{\volt} \)

    \begin{solu}

%        \begin{figure}[H]
%            \centering
%            \includegraphics[scale=.55]{images/imagesUD8ReaccionesRedox/UD8_Problemas/redox_Evau_18.png}
%            %\caption{}
%            %\label{fig:}
%        \end{figure}

        En ambos procesos, el \ch{Fe^2+} se oxida a \ch{Fe^3+} (semirreacción de oxidación, con \(E^0_{ox} = -0,77\ \si{\volt}\)).

        \textbf{a) \ch{Fe^2+ + 1/2 Cl2 -> Fe^3+ + Cl-}}

        \begin{equation*}
            E^0_{cell} = E^0(\ch{Cl2}/\ch{Cl-}) - E^0(\ch{Fe^3+}/\ch{Fe^2+}) = 1,36-0,77 = \SI{0.59}{\volt} > 0
        \end{equation*}

        Como \(E^0_{cell} > 0\), el proceso \textbf{sí es espontáneo}.

        \textbf{b) \ch{Fe^2+ + 1/2 Br2 -> Fe^3+ + Br-}}

        \begin{equation*}
            E^0_{cell} = E^0(\ch{Br2}/\ch{Br-}) - E^0(\ch{Fe^3+}/\ch{Fe^2+}) = 0,54-0,77 = \SI{-0.23}{\volt} < 0
        \end{equation*}

        Como \(E^0_{cell} < 0\), el proceso \textbf{no es espontáneo} (sí lo sería el proceso inverso, la oxidación de \ch{Br-} por \ch{Fe^3+}).

    \end{solu}
\end{exerciseEVAU}

\begin{exerciseEVAU}
    Se construye una pila con los electrodos \(\ch{Mg\pch[2]}\ | \ \ch{Mg}\) y \(\ch{Ag+}\ |\ \ch{Ag}\) en condiciones estándar. a) Escribe las reacciones anódica, catódica y global. b) Escribe la notación esquemática de la pila. c) Indica el sentido del movimiento de las cargas durante el funcionamiento de la pila. d) Calcula su fuerza electromotriz.

    \datos{}: \( E\text{º}(\ch{Mg\pch[2]}\ | \ \ch{Mg}) = \SI{-2.37}{\volt} \); \( E\text{º}(\ch{Ag+}\ | \ \ch{AG}) = \SI{0.78}{\volt} \).

    \begin{solu}

%        \begin{figure}[H]
%            \centering
%            \includegraphics[scale=.55]{images/imagesUD8ReaccionesRedox/UD8_Problemas/redox_Evau_19.png}
%            %\caption{}
%            %\label{fig:}
%        \end{figure}

        Como \(E^0(\ch{Ag+}/\ch{Ag}) = \SI{0.78}{\volt} > E^0(\ch{Mg^2+}/\ch{Mg}) = \SI{-2.37}{\volt}\), la plata actúa como cátodo y el magnesio como ánodo.

        \textbf{a) Reacciones anódica, catódica y global}

        \begin{equation*}
            \text{Ánodo:} \quad \ch{Mg -> Mg^{2+} + 2 e-} \qquad \text{Cátodo:} \quad \ch{2 Ag+ + 2 e- -> 2 Ag}
        \end{equation*}

        \begin{equation*}
            \text{Global:} \quad \ch{Mg + 2 Ag+ -> Mg^{2+} + 2 Ag}
        \end{equation*}

        \textbf{b) Notación esquemática}

        \begin{equation*}
            \ch{Mg}(s) \mid \ch{Mg^{2+}}(aq) \parallel \ch{Ag+}(aq) \mid \ch{Ag}(s)
        \end{equation*}

        \textbf{c) Sentido de circulación de las cargas}

        Los electrones circulan por el circuito externo desde el ánodo (\ch{Mg}) hacia el cátodo (\ch{Ag}). En la disolución, los cationes del puente salino migran hacia el cátodo y los aniones hacia el ánodo, para mantener la electroneutralidad.

        \textbf{d) Fuerza electromotriz}

        \begin{equation*}
            E^0_{pila} = 0,78-(-2,37) = \SI{3.15}{\volt}
        \end{equation*}

    \end{solu}
\end{exerciseEVAU}

\begin{exerciseEVAU}
    Supón una celda voltaica espontánea (pila). Explica razonadamente si las siguientes afirmaciones son ciertas o falsas: a) Los electrones se desplazan del cátodo al ánodo. b) Los electrones atraviesan el puente salino. c) La reducción tiene lugar en el electrodo positivo.

    \begin{solu}

%        \begin{figure}[H]
%            \centering
%            \includegraphics[scale=.55]{images/imagesUD8ReaccionesRedox/UD8_Problemas/redox_Evau_20.png}
%            %\caption{}
%            %\label{fig:}
%        \end{figure}

        \textbf{a) «Los electrones se desplazan del cátodo al ánodo»: FALSO.}

        En una pila voltaica, la oxidación ocurre en el ánodo, que libera electrones, y estos circulan por el circuito externo hacia el cátodo, donde se produce la reducción. Por tanto, los electrones se desplazan \textbf{del ánodo al cátodo}, justo al contrario de lo afirmado.

        \textbf{b) «Los electrones atraviesan el puente salino»: FALSO.}

        Los electrones circulan únicamente por el circuito externo (el conductor metálico que conecta ambos electrodos). El puente salino permite el paso de iones (cationes hacia el cátodo y aniones hacia el ánodo) para mantener la electroneutralidad de ambas disoluciones, pero no es atravesado por electrones.

        \textbf{c) «La reducción tiene lugar en el electrodo positivo»: VERDADERO.}

        En una pila voltaica, el cátodo (donde ocurre la reducción) es, por definición, el electrodo positivo, ya que es hacia donde se dirigen los electrones a través del circuito externo (es el electrodo de mayor potencial).

    \end{solu}
\end{exerciseEVAU}

\begin{exerciseEVAU}
    Se dispone de dos barras metálicas, una de plata, y otra de cinc. También se dispone de las sales de nitrato de estos elementos y cloruro de potasio, material de vidrio adecuado y un voltímetro con conexiones eléctricas. Escribe las reacciones que tienen lugar en el ánodo y en el cátodo de dicha pila indicando qué especie se oxida y cuál se reduce.

    \datos{}: \( E\text{º}(\ch{Ag+}\ | \ \ch{Ag}) = \SI{0.79}{\volt} \); \( E\text{º}(\ch{Zn\pch[2]}\ | \ \ch{Zn}) = \SI{-0.76}{\volt} \)

    \begin{solu}

%        \begin{figure}[H]
%            \centering
%            \includegraphics[scale=.55]{images/imagesUD8ReaccionesRedox/UD8_Problemas/redox_Evau_21.png}
%            %\caption{}
%            %\label{fig:}
%        \end{figure}

        Como \(E^0(\ch{Ag+}/\ch{Ag}) = \SI{0.79}{\volt} > E^0(\ch{Zn^2+}/\ch{Zn}) = \SI{-0.76}{\volt}\), la plata actúa como cátodo (mayor tendencia a reducirse) y el cinc como ánodo.

        \textbf{Ánodo (oxidación): se oxida el \ch{Zn}}

        \begin{equation*}
            \ch{Zn -> Zn^{2+} + 2 e-}
        \end{equation*}

        \textbf{Cátodo (reducción): se reduce el \ch{Ag+}}

        \begin{equation*}
            \ch{2 Ag+ + 2 e- -> 2 Ag}
        \end{equation*}

        La reacción global de la pila es \(\ch{Zn + 2 Ag+ -> Zn^{2+} + 2 Ag}\), con \(E^0_{pila} = 0,79-(-0,76) = \SI{1.55}{\volt}\). El cloruro de potasio (\ch{KCl}) se emplearía para preparar el puente salino que cierra el circuito iónico entre ambas disoluciones.

    \end{solu}
\end{exerciseEVAU}


